% Part: turing-machines
% Chapter: machines-computations
% Section: combining-machines

\documentclass[../../../include/open-logic-section]{subfiles}

\begin{document}

% SEGMENT OLP-0261-S01

\olfileid{tur}{mac}{cmb}
\olsection{டியூரிங் பொறிகளை இணைத்தல்}

\begin{explain}
இதுவரை நாம் பார்த்த டியூரிங் பொறிகளின் எடுத்துக்காட்டுகள் இயல்பில்
மிகவும் எளியவை. ஆனால் உண்மையில், தற்கால நிரலாக்க மொழி எதனாலும்
தீர்க்கக்கூடிய எந்தச் சிக்கலையும் டியூரிங் பொறிகளாலும் தீர்க்கலாம்.
மேலும் சிக்கலான டியூரிங் பொறிகளைக் கட்ட, அவற்றை இணைக்க முடியும் என்று
நமக்கு நாமே உறுதிப்படுத்திக்கொள்வது முக்கியம்; அப்போது செயல்முறையை
எளிய பகுதிகளாகப் பிரித்து, மேலும் சிக்கலான சிக்கல்களைத் தீர்க்கும்
பொறிகளைக் கட்டலாம். ஒரு சிக்கலான சிக்கலை அதன் கூறுப் பகுதிகளாகப்
பிரிப்பதற்கான இயல்பான வழியைக் கண்டறிய முடிந்தால், பல எளிய டியூரிங்
பொறிகளை உருவாக்கி, அவற்றைச் சிக்கலைத் தீர்க்கக்கூடிய ஒரே பொறியாக
இணைப்பதன் மூலம் பல கட்டங்களில் அச்சிக்கலை அணுகலாம். அடுத்த பிரிவில்
நிற்றல் சிக்கலை அணுகும்போது இக்குறிப்பு குறிப்பாக முக்கியமானது.

டியூரிங் பொறிகள் $M = \tuple{Q, \Sigma, q_0, \delta}$ மற்றும்
~$M' = \tuple{Q', \Sigma', q_0', \delta'}$ ஆகியவற்றை எவ்வாறு
இணைப்பது?~$M$ நின்றபிறகு நாடாவின் அமைவுநிலையை இப்போது பொறி~$M'$ இன்
ஓர் இயக்கத்திற்கான உள்ளீட்டு அமைவுநிலையாகப் பயன்படுத்துகிறோம். இதைச்
செய்யும் ஒரே டியூரிங் பொறி $M \frown M'$ ஐப் பெற, பின்வருவனவற்றைச்
செய்க:
\begin{enumerate}
    \item $M$ க்கும்~$M'$ க்கும் பொதுவான நிலைகள் எவையும் இல்லாதவாறு
    ($Q \cap Q' = \emptyset$),~$M'$ இன் எல்லா நிலைகள்~$Q'$ க்கும்
    மறுஎண்ணிடுக (அல்லது மறுபெயரிடுக).
    \item $M \frown M'$ இன் நிலைகள் $Q \cup Q'$.
    \item நாடாவின் குறியெழுத்துத் தொகுதி $\Sigma \cup \Sigma'$.
    \item தொடக்க நிலை~$q_0$.
    \item நிலைமாற்றுச் சார்பு பின்வருமாறு தரப்படும் $\delta''$ என்ற சார்பு:
    \[\delta''(q,\sigma) =
    \begin{cases}
      \delta(q,\sigma) & \text{$q \in Q$ எனில்}\\
      \delta'(q,\sigma) & \text{$q \in Q'$ எனில்}\\
      \tuple{q_0', \sigma, \TMstay} & \text{$q \in Q$ மற்றும்
      $\delta(q,\sigma)$ வரையறுக்கப்படவில்லை எனில்}
    \end{cases}\]
\end{enumerate}
இதனால் கிடைக்கும் பொறி, $q \in Q$ என்ற நிலையில் இருக்கும்போது~$M$ இன்
கட்டளைகளையும்,~$q \in Q'$ என்ற நிலையில் இருக்கும்போது~$M'$ இன்
கட்டளைகளையும் பயன்படுத்துகிறது. $q \in Q$ என்ற நிலையில் இருந்து,
~$M$ இல் நிலைமாற்றம் இல்லாத ஒரு குறி~$\sigma$ ஐ வருடும்போது (அதாவது,
~$M$ நின்றிருக்கும் நிலையில்), அது~$M'$ இன் தொடக்க நிலைக்குள்
செல்கிறது (நாடாவின் உள்ளடக்கத்தையும் தலையின் இருப்பிடத்தையும்
அப்படியே விட்டுவைக்கிறது).

பொறி~$M$ ஒழுங்குக்குட்பட்டதாக இல்லாவிட்டால்,~$M$ நிற்கும்போது
நாடாவின் தலை எங்குள்ளது என்று நமக்குத் தெரியாது. எனவே~$M$ இன்
நிற்றல் அமைவுநிலையில் தலை கட்டம்~$1$ ஐ வருட வேண்டியதில்லை என்பதைக்
கவனிக்க. பொறிகளை இணைக்கும்போது இதை மனத்தில் கொள்வது முக்கியம்.
\end{explain}

\begin{ex}
\emph{பொறிகளை இணைத்தல்:} நீளங்கள் $n$,~$m$ ஆக உள்ள
~$\TMstroke$-களின் இரு தொகுதிகளைக் கொண்ட உள்ளீட்டில் தொடங்கும்போது,
நாடாவில் $2(m+n)$ எண்ணிக்கையிலான $\TMstroke$-களைக் கொண்ட ஒரே
தொகுதியுடன் நிற்கும் ஒரு பொறியை வடிவமைப்போம். இப்பொறியைக் கட்ட,
நமக்கு ஏற்கெனவே பழக்கமான இரு பொறிகளை இணைக்கலாம்: கூட்டல் பொறியையும்
இரட்டிப்பானையும். கூட்டல் பொறிக்கான ஒரு நிலை வரைபடத்தை வரைவதன்
மூலம் தொடங்குகிறோம்.
\[
\begin{tikzpicture}[->,>=stealth',shorten >=1pt,auto,node distance=2.8cm,
                    semithick]
  \tikzstyle{every state}=[fill=none,draw=black,text=black]

  \node[initial,state] (A)              {$q_0$};
  \node[state]         (B) [right of=A] {$q_1$};
  \node[state]         (C) [right of=B] {$q_2$};

  \path (A) edge node {\TMtrans{\TMblank}{\TMstroke}{\TMstay}} (B)
            edge [loop above] node {\TMtrans{\TMstroke}{\TMstroke}{\TMright}} (B)
        (B) edge [loop above] node {\TMtrans{\TMstroke}{\TMstroke}{\TMright}} (B)
            edge node {\TMtrans{\TMblank}{\TMblank}{\TMleft}} (C)
        (C) edge [loop above] node {\TMtrans{\TMstroke}{\TMblank}{\TMstay}} (C);
\end{tikzpicture}
\]
நிலை~$q_2$ இல் நிற்பதற்குப் பதிலாக, வெளியீட்டை இரட்டிப்பதற்காகச்
செயல்பாட்டைத் தொடர விரும்புகிறோம். இரட்டிப்பான் பொறி உள்ளீட்டிலுள்ள
முதல் கோட்டை அழித்து, தனியான வெளியீட்டில் இரு கோடுகளை எழுதுகிறது
என்பதை நினைவுகூர்க. கூட்டல் பொறியின் வெளியீட்டிலுள்ள முதல் கோட்டை
நாடாவின் தலை வாசிப்பதை உறுதிசெய்ய ஒரு கட்டளையைச் சேர்ப்போம்.
\[
\begin{tikzpicture}[->,>=stealth',shorten >=1pt,auto,node distance=2.8cm,
                    semithick]
  \tikzstyle{every state}=[fill=none,draw=black,text=black]

  \node[initial,state] (A)              {$q_0$};
  \node[state]         (B) [right of=A] {$q_1$};
  \node[state]         (C) [right of=B] {$q_2$};
  \node[state]         (D) [below left of=C] {$q_3$};
  \node[state]         (E) [below left of=D] {$q_4$};

  \path (A) edge node {\TMtrans{\TMblank}{\TMstroke}{\TMstay}} (B)
            edge [loop above] node {\TMtrans{\TMstroke}{\TMstroke}{\TMright}} (B)
        (B) edge [loop above] node {\TMtrans{\TMstroke}{\TMstroke}{\TMright}} (B)
            edge node {\TMtrans{\TMblank}{\TMblank}{\TMleft}} (C)
        (C) edge node {\TMtrans{\TMstroke}{\TMblank}{\TMleft}} (D)
        (D) edge [loop left] node {\TMtrans{\TMstroke}{\TMstroke}{\TMleft}} (D)
            edge node[left, xshift=-2mm] {\TMtrans{\TMendtape}{\TMendtape}{\TMright}} (E);
\end{tikzpicture}
\]
உள்ளீட்டை இரட்டிப்பது இப்போது எளிது---நாம் செய்ய வேண்டியதெல்லாம்
இரட்டிப்பான் பொறியை நிலை~$q_4$ உடன் இணைப்பதே. இரட்டிப்பான் பொறியின்
நிலைகள்~$q_0$ குப் பதிலாக~$q_4$ இல் தொடங்குமாறு அவற்றை மறுபெயரிட
வேண்டும்---இவ்வாறு இரு தொடக்க நிலைகள் ஏற்படுவதைத் தவிர்க்கிறோம்.
இறுதி வரைபடம் \olref{fig:combined} இல் உள்ளதுபோல் அமைய வேண்டும்.
\begin{figure}
\[
\begin{tikzpicture}[->,>=stealth',shorten >=1pt,auto,node distance=2.8cm,
                    semithick]
  \tikzstyle{every state}=[fill=none,draw=black,text=black]
  \node[initial,state] (A)              {$q_0$};
  \node[state]         (B) [right of=A] {$q_1$};
  \node[state]         (C) [right of=B] {$q_2$};
  \node[state]         (D) [below left of=C] {$q_3$};
  \node[state]         (E) [below left of=D] {$q_4$};
  \node[state]         (2) [right of=E] {$q_5$};
  \node[state]         (3) [right of=2] {$q_6$};
  \node[state]         (4) [below of=3] {$q_7$};
  \node[state]         (5) [left of=4]  {$q_8$};
  \node[state]         (6) [left of=5]  {$q_9$};

  \path (A) edge node {\TMtrans{\TMblank}{\TMstroke}{\TMstay}} (B)
            edge [loop above] node {\TMtrans{\TMstroke}{\TMstroke}{\TMright}} (B)
        (B) edge [loop above] node {\TMtrans{\TMstroke}{\TMstroke}{\TMright}} (B)
            edge node {\TMtrans{\TMblank}{\TMblank}{\TMleft}} (C)
        (C) edge node {\TMtrans{\TMstroke}{\TMblank}{\TMleft}} (D)
        (D) edge [loop left] node {\TMtrans{\TMstroke}{\TMstroke}{\TMleft}} (D)
            edge node[left, xshift=-2mm] {\TMtrans{\TMendtape}{\TMendtape}{\TMright}} (E)
    (E) edge node {\TMtrans{\TMstroke}{\TMblank}{\TMright}} (2)
    (2) edge [loop above] node {\TMtrans{\TMstroke}{\TMstroke}{\TMright}} (2)
      edge node {\TMtrans{\TMblank}{\TMblank}{\TMright}} (3)
    (3) edge [loop above] node {\TMtrans{\TMstroke}{\TMstroke}{\TMright}} (3)
        edge node {\TMtrans{\TMblank}{\TMstroke}{\TMright}} (4)
    (4) edge [loop below] node {\TMtrans{\TMblank}{\TMstroke}{\TMleft}} (4)
        edge node {\TMtrans{\TMstroke}{\TMstroke}{\TMleft}} (5)
    (5) edge [loop below]  node {\TMtrans{\TMstroke}{\TMstroke}{\TMleft}} (5)
        edge              node {\TMtrans{\TMblank}{\TMblank}{\TMleft}} (6)
    (6) edge [loop below] node {\TMtrans{\TMstroke}{\TMstroke}{\TMleft}} (6)
        edge              node {\TMtrans{\TMblank}{\TMblank}{\TMright}} (E);
\end{tikzpicture}
\]
\caption{கூட்டல் பொறியையும் இரட்டிப்பான் பொறியையும் இணைத்தல்}
\ollabel{fig:combined}
\end{figure}
\end{ex}

\begin{prop}
    $M$, $M'$ ஆகியவை ஒழுங்குக்குட்பட்டவையாக இருந்து, முறையே
    $f\colon \Nat^k \to \Nat$, $f'\colon \Nat \to \Nat$ என்ற
    சார்புகளைக் கணிக்கின்றன எனில், $M \frown M'$ உம்
    ஒழுங்குக்குட்பட்டது; மேலும் அது~$\comp{f}{f'}$ ஐக் கணிக்கிறது.
\end{prop}

\begin{proof}
    $M$ ஒழுங்குக்குட்பட்டது என்பதால், அது
    வெளியீடு~$f(n_1,\dots,n_k) = m$ உடன் நிற்கும்போது, தலை
    கட்டம்~$1$ ஐ வருடுகிறது. இப்போது~$M'$ இன் தொடக்க நிலைக்குள்
    சென்றால், பொறி வெளியீடு $f'(m)$ உடன் மீண்டும் கட்டம்~$1$ ஐ
    வருடியபடி நிற்கும். \olref[dis]{defn:disciplined} இன் மற்ற
    நிபந்தனைகளும் நிறைவேறுகின்றன.
\end{proof}

\begin{prob}
    \cref{tur:mac:dis:prob:disc-succ} இலிருந்து பொறி~$M$ ஐ எடுத்து,
    $M \frown M$ ஐக் கட்டமைப்பதன் மூலம் $f(x) = x+2$ ஐக் கணிக்கும்
    ஒழுங்குக்குட்பட்ட ஒரு டியூரிங் பொறியைத் தருக.
\end{prob}
\end{document}
